\section{Method}\label{sec:method}
\paragraph{Service and tasks.} The agent operates a small in-memory record service with collections
of orders, items and accounts. It reads with \texttt{get\_record} and \texttt{list\_records} and
writes with \texttt{update\_field}, \texttt{transfer}, \texttt{create\_record} and
\texttt{delete\_record}. Write tools echo their effect: an update returns the record after the write
and a transfer returns both balances. Tasks come from six templates (for example, cancel an order and
restock its items, or move money between accounts), each instantiated ten times in each test split and five times in the development split. A task
succeeds only if the agent declares it done, every goal holds in the final state, and nothing else
changed. No language model judges success.

\paragraph{Agent.} Every arm runs the same loop and system prompt. At each step the model replies with
one JSON object, either a tool call with a short thought and a stated confidence (a probability), or a
declaration that the task is done; a missing confidence counts as full confidence. The tool result, or
the tool error, comes back as the next user message. A reply that is not a valid action is returned to
the model as an error; a short run of such replies in a row ends the episode as a \emph{parse
failure}. The agent has at most \RQwenEightMaxSteps{} steps. Sampling settings are fixed per model and
given in the released code.

\paragraph{Faults.} At most one fault is injected per episode. For each (task, fault, seed) a slot
$k$ is drawn and shared by every arm; the fault arms at the $k$-th successful call of its kind and
fires once, on the first eligible call at or after it. Only the slot index is shared: arms change what the agent does, so the
trajectory and whether the fault fires at all can differ between arms (Section~\ref{sec:limitations}). The faults are:
\emph{clean}, which injects nothing;
\emph{loud}, which returns an error;
three silent write faults that return the echo a correct write would return, namely \emph{drop},
which does not apply the write, \emph{wrong}, which writes the value to a different field of the
record, and \emph{partial}, which applies only the debit of a transfer or only the first half of a
new record's fields;
\emph{read}, which corrupts one field of a returned record, with an implausible value (a negative
number or an invalid order status) in a little under half of the episodes and a plausible one otherwise;
and \emph{delayed}, which applies the write but hides it from reads for two further calls. Delayed
exposes false alarms, since there is nothing to repair.

\paragraph{Triggers.} After each tool call a trigger decides whether to flag it. \texttt{never} and
\texttt{always} are the bounds; \texttt{on\_error} fires on tool errors and \texttt{on\_uncertainty}
when the agent's stated confidence is below \RQwenEightThreshold{}. \texttt{on\_contract} sees only
the call and its result and checks the schema, the field types, simple value domains (no negative numbers, a known order
status) and that the echo matches the request.
At the faulted call of a silent write fault the echo is correct, so the contract cannot flag it; it can
flag later reads of a record the fault left with a wrongly typed value or missing fields.
\texttt{on\_readback} computes the effect a write should have had from the harness's own snapshot of
the state before the call and a hand-written effect model per tool, re-reads every affected record, and
fires on any difference. The write faults change the stored state itself, so any read would show them;
read-back's advantage is that it compares what it reads with the expected effect. Its reads go through
a path that the wrong-read fault does not corrupt and that does not count towards fault slots, but
that still sees delayed visibility. This is an idealised read-back: a deployed one would need its own
expected-effect model. \texttt{readback\_contract} fires if either does.

\paragraph{Responses.} A flag injects one fixed message, whatever the trigger:
\begin{quote}\small
NOTICE\@: a check flagged a possible problem with step $k$ (\emph{tool}). Its result may not match the
real state. Decide how to continue: you may re-read state, re-issue the call, or change your plan.
\end{quote}
The agent never sees what the check found. Under \emph{recover} the agent continues with this
message; under \emph{abstain} the episode ends and counts as a failure. All seven triggers are run with
recover, and the contract and read-back triggers also with abstain, giving nine arms.
